\section{第二条线：局地快率怎样连接层长、驱动和全局完成时间}\label{story:global}
局地速率一旦能在多个模型中复现，问题就从“能否快”转向“谁选择这个局地层”。Shay的宽片研究先把漫长等待与短暂快阶段分开；局地解析模型给出速率对上游场和层长的依赖，却把这两个量留给外部解；Sullivan的耗散扫描则说明层长本身也会调整。把这些结果接起来，才看得出显式电阻消去以后，全局问题仍然存在。

受迫边界和磁岛合并进一步改变供磁、背压与完成条件。这里依次阅读局地层选择、不对称与驱动，再读跨模型岛合并和全几何尺寸扫描。它们不是在测同一个数：瞬时电场、固定时间窗平均率、单位X线长度的通量，以及两个岛最终完成合并的时间，必须各自保留定义。

\subsection{嵌入宽电流片：漫长启动与局地快速阶段可以共存}
\citet{Shay2004}研究的初始电流片宽度随系统尺度变化，电子／离子扩散区则嵌在其中。七组二维周期算例中，四个较大算例在岛半宽约 $6.1$--$9.9d_i$ 后进入快速阶段；最大算例最后约3\%的模拟时长完成约90\%的磁通重联与75\%的岛宽增长。这个时间结构是本文的重要发现，不能概括成“从开始到结束都以0.1重联”。

作者实际测量局地上游场 $B_d(t)$：以离子面外电流峰值的25\%位置定义上游边缘，以电子和离子流出速度交点定义下游边缘；启动期另作修正。后期 $V_{\rm out}\propto B_d$，输入速度约为局地Alfvén速度的0.1。若以远处固定磁场归一化，同一时段的率会显著变化。研究没有扫描质量比，故其质量比无关的表述是联系先前工作后的预期，不是本组算例独立证明。

图8还给出一个容易被忽略的现象：小 $\Delta'$ 算例即使已呈开放X形，压力背压仍几乎抵消磁张力，离子喷流及重联仍很慢；较大 $\Delta'$ 的长片则在电子片变薄、Hall贡献增强后快速缩短。因此，单凭某一张X形磁力线图判断快重联并不可靠。默认 $d_e$ 只有两格，另一个高分辨算例还同时降低四阶耗散；这不是严格的独立网格收敛扫描。

\subsection{解析层模型：显式消去电阻，不等于已确定全局速率}
\citet{Malyshkin2008}从X点附近和电子层上游的方程展开得到连接Sweet--Parker与Hall极限的公式；强Hall极限为 $E/(B_m\VA)\sim d_i/L$。但 $L$ 是由上游磁场沿层方向的曲率定义的局地尺度，与 $B_m$ 一起作为输入。原文明确把求解它们留给全局问题。因此，该文解决了“给定局地几何后各项如何平衡”，没有从任意边界与初始电流片推出唯一的0.1。

\citet{Simakov2009}把电子区与离子区作有限区域平衡，以35组岛合并算例检验。其归一化率覆盖约1200倍，模拟／理论比在0.5--2.3之间；强Hall极限得到电子层长度 $w$ 控制的率。这里 $w$ 和局地上游场仍从模拟提取。重要的新结果是：\textbf{细长电子层可以对应对耗散系数不显含的率；开放电子X形也可能处于耗散控制的慢支。}该文的一般过渡公式是渐近匹配插值，不能称为没有近似的全局解析解。

\citet{Malyshkin2009}进一步把电子漂移惯性加入总动量，得到电子半厚约 $d_e$、电子层长度随有效电阻的倒数增长、率随有效电阻正比变化的稳态支。这里“collisionless”指电阻被解释为待动理学决定的有效量，X点仍满足 $E_z=\eta J_z$。它与 \citet{Andres2014} 的无显式耗散、非稳态惯性重联研究的是不同极限；两者并列并不构成逻辑矛盾。

\subsection{电阻从公式中消去，仍有未选定的长度}
\citet{L145}用电阻拖曳同时决定电子排出和耗散层厚，得到$E\sim B_mV_{\!A,m}d_i/L_e$；电阻消去后$L_e$和$B_m$仍是外部输入。
\citet{L137}从给定Harris外场及给定电场构造弱Hall局地解，也未选择全局率，且局部展开出现中心O点等几何问题。
\citet{L160}的有限质量比稳态分支又显示：当电子层厚受到$d_e$限制，$E\sim\eta B_m/(\mu_0d_e)$可显含电阻。
这些研究并不互相取消，而是采用不同闭合、长度约束和稳态假设；“快率与微观项无关”必须附带具体分支条件。

\subsection{电子层长度会依赖耗散，归一化会改变速率的表述}
\citet{Sullivan2009}在等温Hall／惯性双流体岛合并中扫描超电阻。无电子惯性时，电子层长约随 $\eta_H^{1/3}$ 变化，厚度约随 $\eta_H^{4/9}$ 变化；因此不能把局地理论中尚未求出的层长默认为只由边界决定的常数。

更值得学习的是原文图9、10的对照：用局地电子层上游场和密度归一化，速率随超电阻变化；用初始全局量归一化，或直接看原始电场，则速率近乎不变，约在0.1--0.2范围。这不是数据互相矛盾，而是局地归一化量本身随耗散改变。作者把“处理初始岛磁通所需的整体时间”与局地层的归一化几何区分开。

该研究\textbf{没有}重现动理学中长度约 $10d_i$ 的长电子片：其无惯性超电阻扫描最长仍不到一个 $d_i$。惯性扫描所得长度可拟合 $m_e^{3/8}$，但一数量级质量比区间不足以区分 $3/8$ 与 $1/3$。真实质量比下约 $15d_e$ 的长度和约一个 $d_e$ 的厚度均为外推。最低超电阻算例的厚度接近网格，原文也标明不可靠；不能仅凭论文题名中的“延长”断言已解决长电子层问题。

\subsection{非耗散双流体的数值证据与系统尺寸问题}
\citet{Andres2014}的无显式耗散伪谱比较提供了有价值的对照：提高分辨率时，理想MHD与无惯性Hall的伪通量变化趋小，保留电子惯性的模型则趋向非零变化。它比单独展示一张磁岛图更能支持惯性作用。

但总能误差小于 $10^{-8}$ 并不保证所有电子尺度结构准确。真实质量比、Hall参数0.1时，$d_e\simeq0.00233L_0$；$2\pi L_0$ 周期盒的2048网格间距约 $0.00307L_0$。全局通量相近仍不能称作充分解析一个 $d_e$ 宽的局地层。该研究的“惯性必要”限定在标量压力、不可压、无其他非理想项的模型中。

\citet{Andres2016}保持初始几何并扫描Hall参数0.07--0.16，共十组，测得电流层厚度约随 $d_e$ 变化，上游场和层长较弱变化，准稳态平均率约为 $0.11\lambda^{0.98}$。这是有限参数窗中的近线性趋势，不能简写成原归一化下 $\R=\lambda$；论文后一个关系另外把归一化长度和磁场改取局地值。

还须区分数值趋势与解析解释。该预印本的稳态推导有惯性旋度项符号不一致；严格对称、停滞的二维X点若无非对角压力和电阻，稳态对流惯性不能独自支撑有限 $E_z$，必须检查时间惯性或评估位置。报告保留可见的扫描结果，但不把该局地量纲推导当作已闭合的普遍稳态理论。

\subsection{不对称重联：三个特征点分开，标度与绝对数值分开}
\citet{Cassak2008}检验反平行、无导引场的非对称Hall重联。混合Alfvén速度为
\begin{equation}
v_{\rm out}^2\sim\frac{B_1B_2(B_1+B_2)}{\mu_0(\rho_1B_2+\rho_2B_1)}.
\end{equation}
其九组算例的通量率和喷流随不对称参数的变化大体符合该标度；对称基准率实际为约0.0587。值得关注的结构是：磁X点、电子停滞点和离子停滞点可以不同位置；磁岛向弱场一侧鼓起，电流片向强场一侧移动。测流入速度应使用移动X线参考系。

该文另一种离子流出速度定义会改变绝对数值约60\%，但不明显改变参数标度。说明“标度符合”比“所有流速数值相同”弱。计算还采用 $\beta_{\min}=4$、质量比25、惯性项中的常密度近似和四阶扩散。作者把可靠不对称程度限制在约4以内，并指出更大不对称时电子层会失真；不能直接外推到强密度差、导引场存在的磁层顶。

\subsection{受迫重联：同一个“快”字包含不同驱动和结束方式}
\subsubsection{Taylor墙位移：快峰值后仍有漫长慢尾}
\citet{L057}研究最终停止移动的导体墙，给出线性和非线性、MHD和Hall的不同峰值标度。低电阻时，超电阻对X点电场的贡献会超过普通电阻；即使扫描中率对超电阻弱敏感，也不能把实际计算说成只含普通电阻的Hall模型。

任意导引场的四场扩展在$\beta_F=10^{-2}$至$10^2$范围、保持$d_\beta$相同时，得到相似的Hall突发阶段。保持的是$d_\beta=c_\beta d_i$，并非$d_i$固定。典型快峰值约为电阻参照的30倍，并有约8个Alfvén时间的延迟；\textbf{第一轮快爆发只把岛宽增长到最终值约四分之一}，随后仍有慢的Rutherford阶段。因此，“出现快速峰值”不等于“全部受迫磁通迅速完成重联”。\citep{L058}

两文将自身结果与近0.1的说法比较，但驱动、局地场和速率归一化不同。尤其后一文讨论中对上游磁场的线性标度与通常$B_uV_{Au}\propto B_u^2/\sqrt{n_u}$不同，不能只照搬其反驳措辞。更可靠的结论是：有限墙位移问题中的峰率和完成时间需要分别预测。

\subsubsection{体力持续压入：磁通堆积改变速率的分母}
\citet{L076}从较宽、撕裂稳定或边缘稳定的正弦片开始，用局域离子体力驱动入流。较弱驱动可建立平衡；较强驱动先积累上游磁场、压薄片，约到当地$0.7$--$0.8d_i$时出现快速阶段。其具体扫描中的临界驱动幅度在0.005--0.0075之间，这是选定盒子、驱动宽度和闭合的结果。

增加驱动宽度可提高\emph{绝对}电场，却因堆积的$B_u$更强而降低$B_uV_{Au}$归一化后的率。与自发薄片对照时，受迫算例的归一化出流和率可低约50\%。另一个值得单独记录的现象是：\textbf{磁岛宽度几乎饱和，岛内磁场却继续增加，故重联磁通仍增长}。用岛宽作为唯一通量代理会错判重联已经停止。

这篇2005年论文与2008年同题论文研究不同维数，不能合并成一个文件。2005年表I、正文算例交叉引用及部分图的驱动力数字有不一致；报告保留其主要可核实趋势，不为缺失输入自行补出精确对应。

\subsubsection{局域三维驱动：大部分磁通可以绕过重联层}
2008年研究把驱动同时沿第三方向局域化。二维中被驱动区域推来的上游磁通，超过70\%经过重联；相应三维例中约为10\%，其余可以沿第三方向偏转离开。三维局地电子电流还把磁结构向负$z$搬运，最大重联率不再位于驱动中心。持续显著重联的阶段也较短，典型比较约23对38个归一化时间。\citep{F717}

这里三维总磁通由场线追踪和面积积分求得，二维值则按单位第三维长度定义。作者将三维总率除以驱动全宽$2w_z=10d_i$作比较；这是一种有效长度归一化，不是一般三维重联率的唯一定义。典型对照的体力幅度也分别为二维0.15、三维0.2，并非仅改变维数的一对完全相同输入。结论应是\emph{局域第三维提供了磁通旁路，并改变局地层与整体驱动的联系}，而不是三维必然比二维快或慢一个固定倍数。

\subsection{磁岛合并：Hall模型的快峰值可能高估整体合并}
\citet{Ng2015}用同一Fadeev岛合并构型比较MHD、Hall、十矩与PIC。局地Hall层可以很快，但Hall平均率对岛尺度只有较弱依赖，而所比较PIC率约呈 $\lambda^{-0.8}$。Hall结果可高出PIC两到三倍。这里的平均窗是 $0$ 到 $1.5t_A$，不能与另一宽高比或不同平均窗的 $-0.5$ 指数直接混用。

十矩局地闭合通过调节离子松弛尺度可改善小系统压力支撑；大系统却产生PIC中未见的次级磁岛，并缺少相应弹回。\citet{Ng2017}改为非局地闭合后，片宽约 $1.20d_i$，接近PIC的 $1.15d_i$，而Hall约 $0.18d_i$；但平均率尺寸指数约 $-0.45$，仍不同于PIC约 $-0.8$，压力非旋转对称度也不完全匹配。

\citet{Ng2020}的温度梯度闭合进一步改善该尺寸依赖，并出现导引场从0增加到约1时岛合并峰值率先升、到2时再降的结果。不能用“导引场总会抑制重联”替代具体几何的计算。该研究的附录还说明，没有热流的某些稳态十矩停滞点会导致不合理的压力趋零；多演化六个压力分量并不会自动消除闭合难题。

\subsection{Hall模型也能使磁岛弹回，但结果受电子层处理控制}
Bard与Dorelli表明，不含离子动理学的Hall计算也可能在岛合并中出现弹回。在中等网格上，开放的双尺度层与弹回能够同时存在；继续细化后某些小系统不再弹回，而最大系统仍弹回并出现次级岛。这补充了早期“Hall合并过快”的结论：弹回并非动理学独占现象，但不同模型是否给出相同的平均完成时间仍需比较。\citep{L251,Ng2015,Ng2017}

该文以网格代替显式电阻，因此其层形转变并不是固定物理参数下发现了两个连续解。要把“数值耗散减少后的变化”提升为“物理双稳态”，还需要同物理参数、不同历史的持续分支检验；要建立普适合并标度，则需要控制电子层尺度和平均时间窗。

\subsection{导引场有时提高率，也会改变模型之间的差距}
磁岛合并研究\citep{L241}的结果比“导引场抑制重联”更丰富。小系统中从无导引场增加到$B_g=B_0$，PIC与混合模型峰值率分别提高约51\%和53\%，而Hall MHD略降。此时动理学模型中的离子更磁化，离子层变薄，堆积场和喷流增强，Hall与动理学模型的宏观率趋于接近。再增大导引场，则率重新缓慢下降。

无导引场时，PIC离子喷流比Hall结果弱约七倍，离子非旋转对称应力是重要阻力；Hall标量压力不能表达它。导引场较强时虽然宏观喷流相近，压力与应力的各项分解仍不相同。十矩用合适闭合长度可在小系统接近PIC，较大无导引场系统却产生PIC没有的次级岛，说明闭合的可移植性仍是问题。

全局岛合并还会反弹：若片内处理磁通不够快，两个岛先接近，再分开，之后再合并。无导引场PIC平均率随岛尺寸下降较陡；$B_g=B_0$时四种模型的平均率指数接近。大尺度无导引场PIC堆积区出现镜模相容的磁峰和密度凹陷；作者的线性色散检验支持该解释，但它是否进一步\emph{导致}重联减慢尚未证明。

\subsection{强导引场磁绳合并：中间态阻塞，几何改变排岛}
MAST参数的Hall计算展示三类行为：厚于声回旋尺度的稳定碰撞片；中间尺度上迅速产生中心磁岛；薄层形成开放X点及高峰值率。平面算例保留180度对称，中心岛没有优先排出方向，长大后阻塞合并；环形轴对称几何打破这一对称，岛被排出并完成合并。新现象是外部几何可以决定同类局地岛是否终止全局重联。\citep{Stanier2013}

最低超电阻的峰值率较高，却有较慢的早期磁通处理，且开放分离线后来可再闭合、产生偏心岛对。超电阻控制断线同时加热局地电子，改变$\rho_s$；离子黏性也被人为增强。将这一组结果简化为“超电阻越小，重联一直越快”会丢掉文章最重要的非单调和时序信息。

\subsection{全局重联是否存在系统尺寸效应}
\citet{F600}在固定宏观形状的二维Harris片中发现：尺寸增加30倍，PIC的全局平均率约降低2.6倍；Hall扫描增加100倍，率约降低3.4倍。小尺寸Hall结果近似$\ell^{-1/3}$，更大尺寸下降变缓，是否最终平台仍未确定。排气张角变窄伴随率下降；PIC持续产生磁岛，而Hall后期近单X点，Hall中的波列也没有在PIC中同样出现。

这对“跨尺度时一个固定率就足够”提出直接限制，但不构成对所有局地0.1机制的简单否定。初始片厚是否随盒子一起变、比较局地还是远上游场、按物理时间还是相同磁通阶段统计，会改变被回答的问题。第一章压力亏损模型针对局地准稳态；本篇针对整个演化的磁通处理，二者需要在同一诊断下重新比较。该文的平均区间排印矛盾及质量比25使其精确指数仍应审慎使用。

\paragraph{这一线索走到了哪里。}这一问题尚未由“约0.1”解决。局地快层可与慢全局演化、磁岛反弹、驱动磁通旁路和尺寸依赖同时出现。比较Hall与PIC时，最有区分力的是同一供磁过程下的层长、压力阻力、累计重联磁通和完成时间，而不是各自挑一个峰值。后文有限长三维X线还会增加一个几何比例：整条X线中究竟有多长处于活跃状态。
